\section{Nullstellensatz Hilbert}

Nullstellensatz adalah fakta aljabar mendasar yang telah kita gunakan sebelumnya
untuk membenarkan berbagai contoh, dan yang menghubungkan gagasan
$\spec$ dari suatu gelanggang dengan geometri aljabar klasik.

\subsection{Pernyataan dan bukti awal Nullstellensatz}

Nullstellensatz dapat dinyatakan dalam beberapa cara. Mari kita
mulai dengan versi konkret berikut.

\begin{theorem} \label{nullstellensatzoverC}
Semua ideal maksimal dalam gelanggang polinomial $R=\mathbb{C}[x_1, \dots, x_n]$ berasal
dari titik-titik di $\mathbb{C}^n$. Dengan kata lain, jika $\mathfrak{m} \subset R$ adalah
ideal maksimal, maka terdapat $a_1, \dots, a_n \in \mathbb{C}$ sedemikian sehingga
$\mathfrak{m} = (x_1 - a_1, \dots, x_n - a_n)$.
\end{theorem}


Dengan demikian, spektrum maksimal dari $R=\mathbb{C}[x_1, \dots, x_n]$ diidentifikasi dengan
$\mathbb{C}^n$.

Sekarang kita akan mereduksi \rref{nullstellensatzoverC} menjadi pernyataan yang lebih mudah.
Misalkan
$\mathfrak{m}\subset R$ suatu ideal maksimal. Maka terdapat pemetaan
\[ \mathbb{C} \to R \to R/\mathfrak{m}  \]
dan karena itu $R/\mathfrak{m}$ merupakan aljabar-$\mathbb{C}$ terbangkitkan hingga, sebagaimana
$R$. Berdasarkan kemaksimalan, gelanggang $R/\mathfrak{m}$ juga merupakan medan.

Kita ingin menunjukkan bahwa $R/\mathfrak{m}$ merupakan ruang vektor
$\mathbb{C}$ yang terbangkitkan
hingga. Hal ini akan mengimplikasikan bahwa $R/\mathfrak{m}$ integral atas $\mathbb{C}$,
sedangkan $\mathbb{C}$ tidak mempunyai perluasan aljabar sejati. Jadi, jika kita
membuktikannya, pemetaan $\mathbb{C} \to R/\mathfrak{m}$ merupakan suatu
isomorfisme. Jika $a_i \in \mathbb{C}$ ($1 \leq i \leq n$) adalah citra
$x_i $ dalam $R/\mathfrak{m} =
\mathbb{C}$, maka $(x_1 - a_1, \dots, x_n - a_n) \subset
\mathfrak{m}$, sehingga $(x_1 - a_1, \dots, x_n - a_n)=
\mathfrak{m}$.


Akibatnya, Nullstellensatz dalam bentuk ini akan mengikuti dari
pernyataan berikut:

\begin{proposition}
Misalkan $k$ suatu medan dan $L/k$ suatu perluasan medan. Andaikan $L$ merupakan aljabar-$k$
terbangkitkan hingga. Maka $L$ adalah ruang vektor-$k$ berdimensi hingga.
\end{proposition}
Inilah yang akan kita buktikan.

Kita mulai dengan pembuktian mudah dalam kasus khusus berikut.
\begin{lemma}
Andaikan $k$ tak terhitung (misalnya $\mathbb{C}$, kasus awal yang kita minati).
Maka proposisi di atas berlaku.
\end{lemma}
\begin{proof}
Karena $L$ merupakan aljabar-$k$ terbangkitkan hingga, cukup ditunjukkan bahwa $L/k$
bersifat aljabar.
Jika tidak, terdapat $x \in L$ yang tidak aljabar atas $k$. Jadi $x$ tidak memenuhi
polinomial tak nol apa pun.
Kita klaim bahwa unsur-unsur tak terhitung banyaknya $\frac{1}{x-\lambda}, \lambda
\in K$ bebas secara linear
atas $K$. Ini menghasilkan kontradiksi, sebab $L$ merupakan aljabar-$k$
terbangkitkan hingga, sehingga dimensinya atas $k$ paling banyak terhitung. (Perhatikan bahwa
gelanggang polinomial berdimensi terhitung atas $k$, dan $L$ merupakan suatu hasil bagi.)

Mari kita buktikan klaim ini. Andaikan tidak. Maka terdapat ketergantungan linear tak trivial
\[ \sum \frac{c_i}{x - \lambda_i}  = 0, \quad c_i, \lambda_i \in K. \]
Di sini semua $\lambda_j$ berbeda agar relasi tersebut tak trivial. Setelah menyamakan
penyebut, kita memperoleh
\[ \sum_i c_i \prod_{j \neq i } (x- \lambda_j) = 0. \]
Tanpa mengurangi keumuman, $c_1 \neq 0$.
Persamaan ini mula-mula berada di medan $L$. Namun $x$ transenden atas $k$, sehingga kita
dapat memandangnya sebagai relasi dalam gelanggang polinomial.
Dengan memandangnya sebagai relasi dalam gelanggang polinomial, kita melihat bahwa
semua suku dalam jumlah tersebut selain suku $i =1$ habis dibagi oleh $x - \lambda_1$
sebagai polinomial.
Akibatnya, sebagai polinomial dalam peubah tak tentu $x$,
\[ x - \lambda_1 \mid c_1 \prod_{j \neq 1} (x - \lambda_j).  \]
Ini suatu kontradiksi karena semua $\lambda_i$ berbeda.
\end{proof}

Pembuktian ini agak ganjil karena memanfaatkan fakta bahwa $\mathbb{C}$
tak terhitung.
Seharusnya fakta tersebut tidak relevan.

\subsection{Lema normalisasi}

Sekarang mari kita berikan pembuktian yang lebih aljabar.
Kita akan menggunakan fakta berikut yang sangat berguna dalam aljabar komutatif.

\begin{theorem}[Lema normalisasi Noether] Misalkan $k$ suatu medan, dan $R =
k[x_1, \dots, x_n]/\mathfrak{p}$ suatu daerah integral terbangkitkan hingga atas $k$ (dengan
$\mathfrak{p}$ suatu ideal prima dalam gelanggang polinomial).

Maka terdapat subaljabar polinomial $k[y_1, \dots, y_m] \subset R$ sedemikian
sehingga $R$ integral atas $k[y_1, \dots, y_m]$.
\end{theorem}

Kelak akan kita lihat bahwa $m$ adalah \emph{dimensi} dari $R$.

Ada gambaran geometris di balik pernyataan ini. Spektrum $\spec R$ adalah suatu varietas aljabar tak tereduksi
dalam $k^n$ (beserta beberapa titik tambahan), yang berdimensi lebih kecil daripada
$n$ jika $\mathfrak{p} \neq 0$. Maka terdapat suatu \emph{pemetaan hingga} ke $k^m$.
Khususnya, kita dapat memetakan $\spec R \to k^m$ secara surjektif melalui pemetaan integral.
Serat-seratnya bahkan berhingga, sebab keintegralan mengimplikasikan serat berhingga. (Sebenarnya,
kita belum membuktikan hal ini.)

Bagaimana kita menemukan proyeksi hingga seperti itu? Dalam karakteristik
nol, kita cukup mengambil suatu
proyeksi ruang vektor $\mathbb{C}^n \to \mathbb{C}^m$. Untuk proyeksi ``generik''
ke subruang berdimensi sesuai, proyeksi tersebut dapat berperan
sebagai pemetaan hingga kita. Dalam karakteristik $p$, cara ini mungkin tidak berhasil.

\begin{proof}
Pertama, perhatikan bahwa $m$ ditentukan secara tunggal sebagai derajat transendensi
medan hasil bagi $R$ atas $k$.

Di antara peubah-peubah $x_1, \dots, x_n \in R$ (yang, dengan sedikit penyalahgunaan notasi,
kita pandang sebagai unsur $R$), pilih suatu subhimpunan maksimal yang bebas secara aljabar.
Subhimpunan ini tidak mempunyai relasi polinomial tak nol. Khususnya, gelanggang
yang dibangkitkan oleh subhimpunan tersebut hanyalah gelanggang polinomial pada subhimpunan itu.
Kita dapat mempermutasikan peubah-peubah ini dan mengandaikan bahwa
$$\left\{x_1,\dots, x_m\right\}$$ adalah subhimpunan maksimal tersebut. Khususnya, $R$
memuat \emph{gelanggang polinomial} $k[x_1, \dots, x_m]$ dan dibangkitkan oleh
peubah-peubah selebihnya. Namun peubah-peubah selebihnya tidak diadjoin secara bebas
begitu saja.

Strateginya adalah sebagai berikut. Kita akan melakukan berhingga banyaknya perubahan
peubah agar $R$ menjadi integral atas $k[x_1, \dots, x_m]$.

Kasus pokok adalah $m=n-1$. Mari kita tangani kasus ini. Jadi,
\[ R_0 = k[x_1, \dots, x_m]  \subset R = R_0[x_n]/\mathfrak{p}.  \]
Karena $x_n$ tidak bebas secara aljabar, terdapat polinomial tak nol
$f(x_1, \dots, x_m, x_n) \in \mathfrak{p}$.

Kita ingin agar $f$ monik dalam $x_n$. Hal ini akan memberikan keintegralan. Mula-mula,
sifat tersebut belum tentu berlaku. Namun kita akan mengubah sistem koordinat agar demikian.
Pilih $N \gg 0$. Untuk $1 \leq i \leq m$, definisikan
\[ x_i' = x_i + x_n^{N^i}.  \]
Maka persamaannya menjadi
\[ 0 = f(x_1, \dots, x_m, x_n) = f( \left\{x_i' - x_n^{N^i}\right\} , x_n). \]
Sekarang $f(x_1, \dots, x_n, x_{n+1})$ berbentuk suatu jumlah
\[ \sum \lambda_{a_1 \dots b} x_1^{a_1} \dots x_m^{a_m} x_n^{b} , \quad
\lambda_{a_1 \dots b} \in k. \]
Namun $N$ dipilih sangat besar. Mari kita uraikan ekspresi ini dalam $x_i'$ dan
memperhatikan pangkat terbesar $x_n$ yang muncul.
Kita memperoleh bahwa
\[ f(\left\{x_i' - x_n^{N_i}\right\},x_n)
\]
mempunyai pangkat terbesar $x_n$ tepat pada suku yang, dalam ekspresi $f$,
memaksimalkan $a_m$ terlebih dahulu, kemudian
$a_{m-1}, $ dan seterusnya. Eksponen terbesar berbentuk
\[ x_n^{a_m N^m + a_{m-1}N^{m-1} + \dots + b}.	\]
Namun dalam ekspresi
\( f(\left\{x_i' - x_n^{N_i}\right\},x_n)\), kita tidak mungkin memperoleh eksponen $x_n$ selain bentuk-bentuk tersebut. Jika $N$ amat
besar, semua eksponen ini akan saling berbeda.
Khususnya, setiap pangkat $x_n$ muncul tepat satu kali dalam uraian
$f$. Dengan demikian, khususnya $x_n$ integral atas $x_1', \dots, x_n'$.
Karena itu, setiap $x_i$ juga integral.

Jadi kita memperoleh
\begin{quote}
$R$ integral atas $k[x_1', \dots, x_m']$.
\end{quote}

Dengan demikian kita telah membuktikan lema normalisasi dalam kasus kodimensi satu. Bagaimana
dengan kasus umum? Kita mengulangi proses ini.
Andaikan kita mempunyai
\[ k[x_1, \dots, x_m] \subset R.  \]
Misalkan $R'$ subgelanggang $R$ yang dibangkitkan oleh $x_{1}, \dots,x_m, x_{m+1}$. Argumen
yang baru saja diberikan menunjukkan bahwa kita dapat memilih $x_1', \dots, x_m'$ sedemikian sehingga
$R'$ integral atas $k[x_1', \dots, x_m']$, dan $x_i'$
bebas secara aljabar.
Sebenarnya, kita mengetahui bahwa $R' = k[x_1', \dots, x_m', x_{m+1}]$.

Mari kita coba mengulangi argumen sambil meninjau $x_{m+2}$. Misalkan $R'' =
k[x_1', \dots, x_m', x_{m+2}]$ modulo relasi apa pun yang harus dipenuhi oleh $x_{m+2}$.
Jadi, ini merupakan subgelanggang dari $R$. Argumen yang sama menunjukkan bahwa kita dapat
mengubah peubah sedemikian sehingga $x_1'', \dots, x_m''$ bebas secara aljabar
dan $R''$ integral atas $k[x_1'', \dots, x_m'']$. Selain itu,
$k[x_1'', \dots, x_m'', x_{m+2}] = R''$.

Setelah melakukan ini, mari kita berikan argumen untuk $m=n-2$. Dari sini akan tampak
cara menangani kasus umum. Kita klaim bahwa
\begin{quote}
$R$ integral atas $k[x_1'', \dots, x_m'']$.
\end{quote}
Untuk itu, kita perlu memeriksa bahwa $x_{m+1}, x_{m+2}$ integral (sebab unsur-unsur ini
bersama $x''_i$ membangkitkan $R''[x_{m+2}][x_{m+2}]=R$.
Namun berdasarkan konstruksi, $x_{m+2}$ integral atas gelanggang ini. Maka penutupan integral dari
$k[x_1'', \dots, x_{m}'']$ di dalam $R$ memuat
\[ k[x_1'', \dots, x''_m, x_{m+2}] = R''.  \]
Akan tetapi, $R''$ memuat unsur-unsur $x_1', \dots, x_m'$. Berdasarkan konstruksi,
$x_{m+1}$ integral atas $x_1', \dots, x_m'$. Penutupan integral dari
$k[x_1'', \dots, x_{m}'']$ harus memuat $x_{m+2}$. Ini menyelesaikan pembuktian untuk
kasus $m=n-2$. Kasus umum serupa; kita cukup melakukan beberapa perubahan
peubah secara berturut-turut.

\end{proof}
\subsection{Kembali ke Nullstellensatz}

Tinjau suatu aljabar-$k$ terbangkitkan hingga $R$ yang merupakan medan. Kita perlu
menunjukkan bahwa $R$ adalah modul-$k$
hingga. Hal ini akan membuktikan proposisi tersebut.
Perhatikan bahwa $R$ integral atas suatu gelanggang polinomial $k[x_1, \dots, x_m]$ untuk
suatu $m$.
Jika $m > 0$, gelanggang polinomial ini mempunyai lebih dari satu ideal prima,
misalnya $(0)$ dan $(x_1, \dots, x_m)$. Namun ideal-ideal ini harus terangkat menjadi ideal prima di
$R$. Memang, kita telah melihat bahwa untuk setiap perluasan integral, pemetaan
terinduksi pada spektrum bersifat surjektif. Jadi,
\[ \spec R \to \spec k[x_1, \dots, x_m]  \]
surjektif. Jika $R$ suatu medan, berarti $\spec k[x_1, \dots, x_m]$ hanya mempunyai satu
titik, sehingga $m=0$. Jadi $R$ integral atas $k$, dan karenanya aljabar. Karena
$R$ hingga karena terbangkitkan hingga. Dengan demikian terbuktilah satu bentuk
Nullstellensatz.


 Bentuk lain Nullstellensatz, yang
lebih tepat, menyatakan sebagai berikut.

\begin{theorem} \label{gennullstellensatz}
Misalkan $I \subset \mathbb{C}[x_1, \dots, x_n]$. Misalkan $V \subset \mathbb{C}^n$ adalah
subhimpunan $\mathbb{C}^n$ yang didefinisikan oleh ideal $I$ (yakni lokus nol dari
$I$).

Maka $\rad(I)$ tepat merupakan himpunan semua $f$ sedemikian sehingga $f|_V = 0$.
Khususnya,
\[ \rad(I) = \bigcap_{\mathfrak{m} \supset I, \mathfrak{m} \
\mathrm{maksimal}} \mathfrak{m}.  \]
\end{theorem}

Khususnya, terdapat bijeksi antara ideal-ideal radikal dan subhimpunan-subhimpunan aljabar
dari $\mathbb{C}^n$.

Bentuk terakhir teorema tersebut, yang mengikuti dari deskripsi ideal-ideal maksimal
dalam gelanggang polinomial, sangat mirip dengan hasil
\[ \rad(I) = \bigcap_{\mathfrak{p} \supset I, \mathfrak{p} \
\mathrm{prima}} \mathfrak{p},  \]
yang berlaku dalam setiap gelanggang komutatif. Namun hasil umum ini belum tentu
berlaku.

\begin{example}
Irisan semua ideal prima dalam suatu DVR adalah nol, tetapi irisan semua ideal
maksimal tidak nol.
\end{example}

\begin{proof}[Bukti \cref{gennullstellensatz}] 
Sekarang cukup ditunjukkan bahwa untuk setiap ideal prima $\mathfrak{p} \subset \mathbb{C}[x_1,
\dots, x_n]$, berlaku
\[ \mathfrak{p} = \bigcap_{\mathfrak{m} \supset I \ \mathrm{maksimal}}
\mathfrak{m}  \]
sebab setiap ideal radikal merupakan irisan ideal-ideal prima.

Misalkan $R = \mathbb{C}[x_1, \dots,
x_n]/\mathfrak{p}$. Ini suatu daerah integral terbangkitkan hingga atas $\mathbb{C}$. Kita
ingin menunjukkan bahwa irisan semua ideal maksimal dalam $R$ adalah nol. Pernyataan ini
ekuivalen dengan persamaan yang ditampilkan di atas.

Tetapkan $f \in R - \left\{0\right\}$. Misalkan $R'$ pelokalan $R'= R_f$. Maka $R'$ juga suatu
daerah integral yang terbangkitkan hingga atas $\mathbb{C}$.	$R'$ mempunyai suatu ideal maksimal
$\mathfrak{m}$ (yang mula-mula mungkin saja nol). Jika kita
meninjau pemetaan $R' \to R'/\mathfrak{m}$, diperoleh pemetaan ke suatu medan yang
terbangkitkan hingga
atas $\mathbb{C}$, sehingga medan itu sama dengan $\mathbb{C}$.
Pemetaan komposit
\[ R \to  R' \to R'/\mathfrak{m}   \]
ditentukan oleh suatu $n$-tupel bilangan kompleks, yakni oleh suatu titik di
$\mathbb{C}^n$ yang bahkan berada di $V$, sebab pemetaan tersebut berasal dari $R$. Titik ini bersesuaian
dengan suatu ideal maksimal dalam $R$.
Berdasarkan konstruksi, ideal maksimal ini tidak memuat $f$.
\end{proof}

\begin{exercise} 
Buktikan hasil berikut, yang dikenal sebagai ``Lema Zariski'' (dan dengan mudah mengimplikasikan
Nullstellensatz): jika $k$ suatu medan dan $k'$ suatu perluasan medan dari $k$ yang
merupakan \emph{aljabar}-$k$ terbangkitkan hingga, maka $k'$ aljabar hingga atas
$k$. Gunakan argumen McCabe berikut (dalam \cite{Mc76}):
\begin{enumerate}
\item $k'$ memuat suatu subgelanggang $S$ berbentuk $S= k[ x_1, \dots, x_t]$, dengan
$x_1, \dots, x_t$ bebas secara aljabar atas $k$, dan $k'$ aljabar
atas medan hasil bagi $S$ (yang merupakan suatu gelanggang polinomial). 
\item Jika $k'$ tidak aljabar atas $k$, maka $S \neq k$ bukan suatu medan.
\item Tunjukkan bahwa terdapat $y \in S$ sedemikian sehingga $k'$ integral atas $S_y$.
Simpulkan bahwa $S_y$ suatu medan.
\item Karena $\spec (S_y) = \left\{0\right\}$, buktikan bahwa $y$ berada dalam setiap
ideal prima tak nol dari $\spec S$. Simpulkan bahwa $1+y \in k$, sehingga $S$ suatu
medan---kontradiksi.
\end{enumerate}
\end{exercise} 

\subsection{Sedikit geometri aljabar afin}

Dalam pembahasan berikut, misalkan $k$ tertutup aljabar dan $A$ suatu aljabar-$k$ terbangkitkan hingga. Ingat bahwa $\Specm A$ menyatakan himpunan ideal maksimal dalam $A$. Tinjau struktur aljabar-$k$ alami pada $\mathrm{Funct}(\Specm A, k)$. Kita mempunyai pemetaan
$$A \rightarrow \mathrm{Funct}(\Specm A, k)$$
yang diperoleh dari Nullstellensatz Lemah sebagai berikut. Ideal-ideal maksimal $\mathfrak{m}\subset A$ berkorespondensi bijektif dengan pemetaan $\varphi_\mathfrak{m}:A\rightarrow k$ yang memenuhi $\ker(\varphi_\mathfrak{m})=\mathfrak{m}$; jadi kita definisikan $a\longmapsto [\mathfrak{m}\longmapsto \varphi_\mathfrak{m}(a)]$. Jika $A$ tereduksi, pemetaan ini injektif: jika $a\in A$ dipetakan ke fungsi nol, maka $a\in \cap\, \mathfrak{m}$ $\rightarrow$ $a$ nilpoten $\rightarrow$ $a=0$.\\

\begin{definition} Suatu fungsi $f\in \mathrm{Funct}(\Specm A,k)$ disebut {\bf aljabar} jika berada dalam citra $A$ di bawah pemetaan di atas. (Istilah lain untuk sifat ini adalah {\bf polinomial} dan {\bf reguler}.) \end{definition}

Misalkan $A$ dan $B$ aljabar-$k$ terbangkitkan hingga, dan $\phi:A\rightarrow B$ suatu homomorfisme. Pemetaan ini menghasilkan pemetaan $\Phi:\Specm B\rightarrow \Specm A$ dengan mengambil prapeta.

\begin{definition} Suatu pemetaan $\Phi:\Specm B\rightarrow \Specm A$ disebut {\bf aljabar} jika berasal dari suatu homomorfisme $\phi$ seperti di atas.\end{definition}

Untuk menunjukkan hubungan antara kedua definisi ini, kita mempunyai proposisi berikut.

\begin{proposition} Suatu pemetaan $\Phi:\Specm B\rightarrow \Specm A$ bersifat aljabar jika dan hanya jika untuk setiap fungsi aljabar $f\in \mathrm{Funct}(\Specm A,k)$, tarikan balik $f\circ \Phi\in \mathrm{Funct}(\Specm B,k)$ bersifat aljabar.\end{proposition}

\begin{proof}
Andaikan $\Phi$ aljabar. Cukup diperiksa bahwa diagram berikut komutatif:
\[ 
\xymatrix{
\mathrm{Funct}(\Specm A,k) \ar[r]^{-\circ\Phi} & \mathrm{Funct}(\Specm B,k) \\
A \ar[u] \ar[r]_{\phi} & B \ar[u]\\
}
\]
dengan $\phi:A\rightarrow B$ pemetaan yang menimbulkan $\Phi$.

[$\Leftarrow$] Andaikan untuk setiap fungsi aljabar $f\in \mathrm{Funct}(\Specm A,k)$, tarikan balik $f\circ\Phi$ bersifat aljabar. Maka diperoleh pemetaan terinduksi dengan menelusuri diagram berlawanan arah jarum jam:
$$
\xymatrix{
\mathrm{Funct}(\Specm A,k) \ar[r]^{-\circ\Phi} & \mathrm{Funct}(\Specm B,k) \\
A \ar[u] \ar@{-->}[r]_{\phi} & B \ar[u]\\
}
$$
Dari $\phi$, kita dapat membangun pemetaan $\Phi':\Specm B \rightarrow \Specm A$ yang diberikan oleh $\Phi'(\mathfrak{m})=\phi^{-1}(\mathfrak{m})$. Kita klaim bahwa $\Phi=\Phi'$. Jika tidak, untuk suatu $\mathfrak{m}\in \Specm B$ berlaku $\Phi(\mathfrak{m})\neq \Phi'(\mathfrak{m})$. Berdasarkan definisi, untuk setiap fungsi aljabar $f\in \mathrm{Funct}(\Specm A,k)$, berlaku $f\circ\Phi=f\circ\Phi'$; untuk memperoleh kontradiksi, kita buktikan lema berikut:\\
Diberikan dua titik berbeda dalam $\Specm A=V(I)\subset k^n$, terdapat suatu
fungsi aljabar $f$ yang memisahkan keduanya. Ini langsung terlihat ketika kita menyadari bahwa setiap
fungsi polinomial bersifat aljabar, dan polinomial-polinomial semacam itu memisahkan titik.
\end{proof}
	
